% Source: https://artofproblemsolving.com/community/c3h1622208_trigonometry
% Stable identity: aops-topic-1622208-q01
Let the rectangle be $ABCD$ with width $AD=w$ and length $AB=w+10$, and let the diagonals meet at $O$. The diagonals of a rectangle are equal and bisect each other, so $OA=OD$ and triangle $AOD$ is isosceles. The diagonals make two pairs of angles, $78^\circ$ and $180^\circ-78^\circ=102^\circ$; the smaller angle faces the shorter side, so $\angle AOD=78^\circ$.

\begin{asy}
usepackage("amsmath");
size(6.5cm);
defaultpen(fontsize(10pt));
real t=tan(39*pi/180); real w=t/(1-t); real L=w+1;
pair A=(0,0),B=(L,0),C=(L,w),D=(0,w),O=(L/2,w/2);
draw(A--B--C--D--cycle,linewidth(0.9));draw(A--C);draw(B--D);
draw(arc(O,0.8,141,219),red);
label("$78^\circ$",O+(-1.3,0),red);
label("$A$",A,SW);label("$B$",B,SE);label("$C$",C,NE);label("$D$",D,NW);label("$O$",O,2.2*E);
label("$w$",(0,w/2),W);label("$w+10$",(L/2,0),S);
label("\mbox{}",(min(currentpicture,user=true)+max(currentpicture,user=true))/2,black); // keeps the text after the figure black when LaTeX typesets the labels inline
\end{asy}

The perpendicular from $O$ to $AD$ (not drawn) bisects $AD$ and the angle $AOD$. It has length $\dfrac{AB}{2}=\dfrac{w+10}{2}$, so in the right triangle with angle $39^\circ$ at $O$,
\[
\tan39^\circ=\dfrac{\dfrac {w}{2}}{\dfrac{w+10}{2}}=\dfrac{w}{w+10}.
\]
Solving, $w=(w+10)\tan39^\circ$, so
\[
w=\dfrac{10\tan39^\circ}{1-\tan39^\circ}\approx\dfrac{8.0978}{0.19022}\approx42.572 ,\qquad w+10\approx52.572 .
\]
The perimeter is
\[
2(w+w+10)=4w+20\approx\boxed{190.3\text{ units}}.
\]
